\begin{frame}{Caso integrador: Empresa Peruana ABC S.A.C.}
\small
\begin{columns}[T,totalwidth=\textwidth]
\column{0.55\textwidth}
\casehead{Caso académico ficticio elaborado con fines educativos.}

\vspace{0.4em}
\begin{tabular}{lr}
\toprule
Dato & Importe \\
\midrule
Activo corriente & \soles{540,000.00}\\
Inventarios & \soles{180,000.00}\\
Activo total & \soles{1,200,000.00}\\
Pasivo corriente & \soles{300,000.00}\\
Pasivo total & \soles{570,000.00}\\
Patrimonio & \soles{630,000.00}\\
Ventas netas & \soles{1,500,000.00}\\
Utilidad neta & \soles{126,900.00}\\
\bottomrule
\end{tabular}

\column{0.40\textwidth}
\centering
\begin{tikzpicture}[node distance=0.45cm, every node/.style={draw, rounded corners=2pt, align=center, font=\scriptsize, minimum width=3.1cm}]
\node (a) {Estado de Situación\\Financiera};
\node[below=of a] (b) {Estado de\\Resultados};
\node[below=of b, fill=unacblue!8] (c) {Ratios\\financieros};
\draw[-{Stealth}] (a) -- (b);
\draw[-{Stealth}] (b) -- (c);
\end{tikzpicture}

\vspace{0.8em}
\scriptsize Esta diapositiva presenta los datos base; los cálculos se desarrollan en la siguiente.
\end{columns}
\end{frame}

\begin{frame}{Caso integrador: operaciones y respuestas}
\scriptsize
\begin{tabularx}{\textwidth}{p{0.24\textwidth}X p{0.16\textwidth}}
\toprule
Ratio & Operación & Resultado \\
\midrule
Razón corriente & \soles{540,000.00} / \soles{300,000.00} & 1.80 \\
Prueba ácida & (\soles{540,000.00} - \soles{180,000.00}) / \soles{300,000.00} & 1.20 \\
Endeudamiento & \soles{570,000.00} / \soles{1,200,000.00} x 100 & 47.50 \pct \\
Margen neto & \soles{126,900.00} / \soles{1,500,000.00} x 100 & 8.46 \pct \\
ROA simple & \soles{126,900.00} / \soles{1,200,000.00} x 100 & 10.58 \pct \\
ROE simple & \soles{126,900.00} / \soles{630,000.00} x 100 & 20.14 \pct \\
\bottomrule
\end{tabularx}

\vspace{0.5em}
\begin{block}{Aclaración}
El ROA y el ROE se calculan con saldos finales debido a que el caso no proporciona valores iniciales. En un análisis más riguroso se emplearían activos y patrimonio promedio.
\end{block}
\end{frame}

\begin{frame}{Diagnóstico y decisiones}
\small
\begin{columns}[T,totalwidth=\textwidth]
\column{0.50\textwidth}
\begin{tabular}{p{0.31\textwidth}p{0.60\textwidth}}
\toprule
Área & Diagnóstico \\
\midrule
Liquidez & 1.80 baja a 1.20 al excluir inventarios. \\
Solvencia & 47.50 \pct{} de activos financiados por terceros. \\
Rentabilidad & Margen neto de 8.46 \pct. \\
Propietarios & ROE simple de 20.14 \pct, ligado al endeudamiento. \\
\bottomrule
\end{tabular}

\column{0.47\textwidth}
\casehead{Decisiones sugeridas}
\begin{itemize}
  \item Fortalecer la cobranza.
  \item Controlar inventarios.
  \item Revisar costo financiero.
  \item Comparar márgenes anteriores.
  \item Implementar seguimiento mensual.
\end{itemize}
\end{columns}
\vspace{0.3em}
\centering
\begin{tikzpicture}[node distance=0.28cm, every node/.style={font=\scriptsize, draw, rounded corners=2pt, align=center}]
\node (a) {Datos};
\node[right=of a] (b) {Ratios};
\node[right=of b] (c) {Diagnóstico};
\node[right=of c] (d) {Decisión};
\node[right=of d] (e) {Seguimiento};
\draw[-{Stealth}] (a)--(b);\draw[-{Stealth}] (b)--(c);\draw[-{Stealth}] (c)--(d);\draw[-{Stealth}] (d)--(e);
\end{tikzpicture}

\scriptsize Los ratios pueden automatizarse mediante ERP, bases de datos, Excel, Python o Power BI, pero la decisión requiere interpretación contable.
\end{frame}
